% Matrix inverse: one LU factorization, then one solve per column of the
% identity.  Cropped to its own bounding box so it can be dropped onto a
% Xournal++ page as an image.  Same style as gauss-pseudocode.tex.
%
%   latexmk -pdf inverse-pseudocode.tex
\documentclass[border=6pt,varwidth=11cm]{standalone}

\usepackage{algorithm2e}
\usepackage{amsmath}

\RestyleAlgo{plain}          % no float, no ruled box
\DontPrintSemicolon
\SetKwComment{Comment}{// }{}
\SetKwInOut{Input}{input}
\SetKwInOut{Output}{output}

\begin{document}
\begin{minipage}{\linewidth}
\begin{algorithm}[H]
  \textbf{Matrix Inverse via LU factorization}
  \BlankLine
  \Input{$n \times n$ matrix $[A]$}
  \Output{$[A]^{-1}$}
  \BlankLine
  $[L], [U] \leftarrow \texttt{lu\_decompose}([A])$\;
  \For{$k \leftarrow 1$ \KwTo $n$}{
    \tcp{$\{e_k\}$ is the $k$th column of $[I]$:}
    \tcp{all zeros, with a 1 in row $k$}
    $A^{-1}[:,k] \leftarrow \texttt{lu\_solve}([L], [U], \{e_k\})$\;
  }
\end{algorithm}
\end{minipage}
\end{document}
